(a) Where $x_0x_1x_2\ne0$, applying $\varphi$ twice gives $(x_0^2x_1x_2:x_0x_1^2x_2:x_0x_1x_2^2)=(x_0:x_1:x_2)$. Thus $\varphi$ is birational and is its own inverse.

(b) Take $U=V=D_+(x_0x_1x_2)$. The displayed formula and its identical inverse are regular there, so $\varphi:U\to V$ is an isomorphism.

(c) Both maximal domains are $\mathbb P^2$ minus the three coordinate vertices. The formula defines a morphism there. At $(0:0:1)$, restriction to $x_0=0$ is constantly $(1:0:0)$ away from the vertex, whereas restriction to $x_1=0$ is constantly $(0:1:0)$. Thus no continuous extension is possible. The other vertices are excluded by symmetry. Each punctured coordinate line is contracted to the opposite vertex.
